\documentclass[12pt,letterpaper]{article}
\usepackage{graphicx,textcomp}
\usepackage{natbib}
\usepackage{setspace}
\usepackage{fullpage}
\usepackage{color}
\usepackage[reqno]{amsmath}
\usepackage{amsthm}
\usepackage{fancyvrb}
\usepackage{amssymb,enumerate}
\usepackage[all]{xy}
\usepackage{endnotes}
\usepackage{lscape}
\newtheorem{com}{Comment}
\usepackage{float}
\usepackage{hyperref}
\newtheorem{lem} {Lemma}
\newtheorem{prop}{Proposition}
\newtheorem{thm}{Theorem}
\newtheorem{defn}{Definition}
\newtheorem{cor}{Corollary}
\newtheorem{obs}{Observation}
\usepackage[compact]{titlesec}
\usepackage{dcolumn}
\usepackage{tikz}
\usetikzlibrary{arrows}
\usepackage{multirow}
\usepackage{xcolor}
\newcolumntype{.}{D{.}{.}{-1}}
\newcolumntype{d}[1]{D{.}{.}{#1}}
\definecolor{light-gray}{gray}{0.65}
\usepackage{url}
\usepackage{listings}
\usepackage{color}

\definecolor{codegreen}{rgb}{0,0.6,0}
\definecolor{codegray}{rgb}{0.5,0.5,0.5}
\definecolor{codepurple}{rgb}{0.58,0,0.82}
\definecolor{backcolour}{rgb}{0.95,0.95,0.92}

\lstdefinestyle{mystyle}{
	backgroundcolor=\color{backcolour},   
	commentstyle=\color{codegreen},
	keywordstyle=\color{magenta},
	numberstyle=\tiny\color{codegray},
	stringstyle=\color{codepurple},
	basicstyle=\footnotesize,
	breakatwhitespace=false,         
	breaklines=true,                 
	captionpos=b,                    
	keepspaces=true,                 
	numbers=left,                    
	numbersep=5pt,                  
	showspaces=false,                
	showstringspaces=false,
	showtabs=false,                  
	tabsize=2
}
\lstset{style=mystyle}
\newcommand{\Sref}[1]{Section~\ref{#1}}
\newtheorem{hyp}{Hypothesis}

\title{Problem Set 3}
\date{Due: March 25, 2026}
\author{Applied Stats II}


\begin{document}
	\maketitle
	\section*{Instructions}
	\begin{itemize}
	\item Please show your work! You may lose points by simply writing in the answer. If the problem requires you to execute commands in \texttt{R}, please include the code you used to get your answers. Please also include the \texttt{.R} file that contains your code. If you are not sure if work needs to be shown for a particular problem, please ask.
\item Your homework should be submitted electronically on GitHub in \texttt{.pdf} form.
\item This problem set is due before 23:59 on Wednesday March 25, 2026. No late assignments will be accepted.
	\end{itemize}

	\vspace{.25cm}
\section*{Question 1}
\vspace{.25cm}
\noindent We are interested in how governments' management of public resources impacts economic prosperity. Our data come from \href{https://www.researchgate.net/profile/Adam_Przeworski/publication/240357392_Classifying_Political_Regimes/links/0deec532194849aefa000000/Classifying-Political-Regimes.pdf}{Alvarez, Cheibub, Limongi, and Przeworski (1996)} and is labelled \texttt{gdpChange.csv} on GitHub. The dataset covers 135 countries observed between 1950 or the year of independence or the first year forwhich data on economic growth are available ("entry year"), and 1990 or the last year for which data on economic growth are available ("exit year"). The unit of analysis is a particular country during a particular year, for a total $>$ 3,500 observations. 

\begin{itemize}
	\item
	Response variable: 
	\begin{itemize}
		\item \texttt{GDPWdiff}: Difference in GDP between year $t$ and $t-1$. Possible categories include: "positive", "negative", or "no change"
	\end{itemize}
	\item
	Explanatory variables: 
	\begin{itemize}
		\item
		\texttt{REG}: 1=Democracy; 0=Non-Democracy
		\item
		\texttt{OIL}: 1=if the average ratio of fuel exports to total exports in 1984-86 exceeded 50\%; 0= otherwise
	\end{itemize}
	
\end{itemize}
\newpage
\noindent Please answer the following questions:

\begin{enumerate}
	\item Construct and interpret an unordered multinomial logit with \texttt{GDPWdiff} as the output and "no change" as the reference category, including the estimated cutoff points and coefficients.
	\item Construct and interpret an ordered multinomial logit with \texttt{GDPWdiff} as the outcome variable, including the estimated cutoff points and coefficients. \newpage
	
	\section{Answers}
	\subsection{Code 1A}
	
	\tiny \begin{BVerbatim}
		gdp_q1 <- na.omit(gdp_data[, c("GDPWdiff", "REG", "OIL")])
		
		gdp_q1$GDPWdiff_cat <- case_when(
		gdp_q1$GDPWdiff < 0  ~ "negative",
		gdp_q1$GDPWdiff > 0  ~ "positive",
		gdp_q1$GDPWdiff == 0 ~ "no change"
		)
		
		table(gdp_q1$GDPWdiff_cat)
		
		gdp_q1$GDPWdiff_unordered <- relevel(
		factor(gdp_q1$GDPWdiff_cat, levels = c("no change", "negative", "positive")),
		ref = "no change"
		)
		
		mnl_mod <- multinom(GDPWdiff_unordered ~ REG + OIL, data = gdp_q1, trace = FALSE)
		summary(mnl_mod)
		
		mnl_z <- summary(mnl_mod)$coefficients / summary(mnl_mod)$standard.errors
		mnl_p <- (1 - pnorm(abs(mnl_z), 0, 1)) * 2
		
		mnl_z
		mnl_p
		
		exp(coef(mnl_mod))

	\end{BVerbatim}
	\normalsize
	
	\subsection{A: Unordered Multinomial logit model}
	\noindent I start by loading the data and checking its structure using \texttt{str()}. When I check the structure of \texttt{GDPWdiff} specifically I see that it is an integer and numbers can be positive, negative and potentially also zero if there is zero `raw difference' in GDP (i.e. it remained equal to the previous year).\newline
	
\noindent Next,	I select the relevant variables as asked in the problem set and drop the missing observations. I then create a categorical variable instead, to observe the numerical `raw difference' in GDP change compared to the previous year. This means that negative is a decrease in GDP, positive implies an increase in GDP and equal to zero means GDP that year is equal to what it was in the previous year. \newline

\noindent Next, I check the distribution of change in GDP in a given year using \texttt{table()}, which shows that in total there are 1105 observations where the GDP change was negative meaning that it decreased. Furthermore, there are 16 zeros and there are 2600 observations where GDP increased, covering the period between 1950 or the year of independence or the first year for which data on economic growth are available (``entry year''), and 1990 or the last year for which data on economic growth are available (``exit year'').\newline
	
	\noindent Next, I apply an unordered multinomial logit model. In this model, I assume no order between the categories; positive, negative and no change are simply three separate options. The assignment asks for the `no change' category as reference/baseline category, meaning all coefficients will be interpreted relative to `no change'. I use \texttt{multinom()} from the \texttt{nnet} package.\newline
	
	\noindent Also, since \texttt{multinom()} I calculate the p-values manually by calculating the z-scores. The z-score tells us how many standard errors the coefficient is away from zero. And the further from zero, the less likely it is due to coincidence. The p-value is the probability of observing a z-value this extreme if H0 (no effect) were true. I multiply by 2 because it is a two-sided test.\newline

	
	\begin{table}[H]
		\centering
		\caption{Unordered Multinomial Logit Model (Reference: No Change)}
		\label{tab:mnl}
		\begin{tabular}{lcccc}
			\hline\hline
			& Intercept       & Democracy (REG)             & Oil dependence (OIL)    \\
			\hline
			\multicolumn{4}{l}{\textit{Coefficients}} \\
			Negative          & 3.8054$^{***}$  & 1.3793$^{*}$    & 4.7840 \\
			Positive          & 4.5338$^{***}$  & 1.7690$^{**}$   & 4.5763 \\
			\hline
			\multicolumn{4}{l}{\textit{Model fit}} \\
			Residual Deviance & \multicolumn{3}{c}{4678.77} \\
			AIC               & \multicolumn{3}{c}{4690.77} \\
			\hline\hline
			\multicolumn{4}{l}{\footnotesize $^{*}$p$<$0.1; $^{**}$p$<$0.05; $^{***}$p$<$0.01} \\
			\multicolumn{4}{l}{\footnotesize Note: Outcome variable is GDPWdiff (negative/positive vs.\ no change).} \\
			\multicolumn{4}{l}{\footnotesize REG: 1 = Democracy; OIL: 1 = fuel exports $>$ 50\%.} \\
		\end{tabular}
	\end{table}
\end{enumerate}

\noindent The results of the unordered multinomial logit model are presented in Table 1. 
Looking at the odds ratios, I can interpret the coefficients as follows.\newline

\noindent We have two independent  variables \texttt{textREG} democracy  and \texttt{textOIL}. The latter takes 1 if the average ratio of fuel exports to total exports in the period 1984-1986 exceeded 50 percent of the total exports. The aforementioned takes 1 if the country is a democracy. \newline

\noindent The results imply that in a democracy, the log-odds of a negative or positive GDP change versus no absolute change are higher compared to the reference category non-democracy, holding everything else constant. \newline


\noindent However, while the effect of democracy on positive growth is statistically significant (p = 0.0210), the effect of demorcacy on negative growth is only marginally significant (p $<$ 0.0727). Furthermore, the independent variable \texttt{textOIL} does not have siginifant effects (p $<$ 0.4871792 for negative \& p $<$  0.5062612 for positive). This could be due to the fact that we have so little observations in our baseline category (only 16).


\subsection{Code 1B}

\tiny \begin{BVerbatim}
	gdp_q1$GDPWdiff_ordered <- ordered(
	gdp_q1$GDPWdiff_cat,
	levels = c("negative", "no change", "positive")
	)
	
	ologit_mod <- polr(GDPWdiff_ordered ~ REG + OIL, data = gdp_q1,
	method = "logistic", Hess = TRUE)
	summary(ologit_mod)
	
	ctable <- coef(summary(ologit_mod))
	p <- pnorm(abs(ctable[, "t value"]), lower.tail = FALSE) * 2
	(ctable <- cbind(ctable, "p value" = p))
	
	exp(coef(ologit_mod))
	
	ologit_mod$zeta
\end{BVerbatim}
\normalsize \newline
\subsection{B: Ordered Multinomial Logit Model}
\noindent For the ordered logit model, unlike in part A, I now assume that there is a natural order between the categories: negative $<$ no change $<$ positive.  I create an ordered factor to reflect this ordering and apply \texttt{polr()} from the \texttt{MASS} package loaded earlier in the code with \texttt{Hess = TRUE} Later, I calculate the p-values manually following the same approach as in part A.\newline

\begin{table}[H]
	\centering
	\caption{Ordered Multinomial Logit Model (negative $<$ no change $<$ positive)}
	\label{tab:ologit}
	\begin{tabular}{lccc}
		\hline\hline
		& Value   & Std. Error & p value \\
		\hline
		\multicolumn{4}{l}{\textit{Coefficients}} \\
		REG                & 0.3985$^{***}$ & 0.0752 & 0.0000 \\
		OIL                & -0.1987$^{*}$  & 0.1157 & 0.0859 \\
		\hline
		\multicolumn{4}{l}{\textit{Cutoff points}} \\
		negative$|$no change & -0.7312$^{***}$ & 0.0476 & 0.0000 \\
		no change$|$positive & -0.7105$^{***}$ & 0.0475 & 0.0000 \\
		\hline
		\multicolumn{4}{l}{\textit{Model fit}} \\
		Residual Deviance  & \multicolumn{3}{c}{4687.689} \\
		AIC                & \multicolumn{3}{c}{4695.689} \\
		\hline\hline
		\multicolumn{4}{l}{\footnotesize $^{*}$p$<$0.1; $^{**}$p$<$0.05; $^{***}$p$<$0.01} \\
		\multicolumn{4}{l}{\footnotesize REG: 1 = Democracy; OIL: 1 = fuel exports $>$ 50\%.} \\
	\end{tabular}
\end{table}

\noindent The results of the ordered multinomial logit model are presented in Table 2. \newline

\noindent For democracy \texttt{REG}, the coefficient is statistically significant 
(p $<$ 0.001, p = 1.16e-07), with an odds ratio of 1.4895639. This means that in a democracy, the cumulative odds of being in a higher ranked category are 49\% higher compared to a non-democracy, holding everything else constant. In simple terms, democracies are more likely to experience positive GDP growth than non-democracies.\newline

\noindent For oil dependence \texttt{OIL}, the coefficient is only marginally significant (p $<$ 0.10, p = 0.086).\newline

\noindent The logit model also provides two cutoff points, where the model draws the estimated boundary between the ordered categories. These values are very close to each other (around -73), which again could be due to the very small size of the middle category (only 16 observations). \newline

\section*{Question 2} 
\vspace{.25cm}

\noindent Consider the data set \texttt{MexicoMuniData.csv}, which includes municipal-level information from Mexico. The outcome of interest is the number of times the winning PAN presidential candidate in 2006 (\texttt{PAN.visits.06}) visited a district leading up to the 2009 federal elections, which is a count. Our main predictor of interest is whether the district was highly contested, or whether it was not (the PAN or their opponents have electoral security) in the previous federal elections during 2000 (\texttt{competitive.district}), which is binary (1=close/swing district, 0="safe seat"). We also include \texttt{marginality.06} (a measure of poverty) and \texttt{PAN.governor.06} (a dummy for whether the state has a PAN-affiliated governor) as additional control variables. 

\begin{enumerate}
	\item [(a)]
	Run a Poisson regression because the outcome is a count variable. Is there evidence that PAN presidential candidates visit swing districts more? Provide a test statistic and p-value.

	\item [(b)]
	Interpret the \texttt{marginality.06} and \texttt{PAN.governor.06} coefficients.
	
	\item [(c)]
	Provide the estimated mean number of visits from the winning PAN presidential candidate for a hypothetical district that was competitive (\texttt{competitive.district}=1), had an average poverty level (\texttt{marginality.06} = 0), and a PAN governor (\texttt{PAN.governor.06}=1).
	
\end{enumerate}


\newpage
\section*{Question 2}

\subsection{Answers Question 2}

\noindent The outcome variable \texttt{PAN.visits.06} is a count, it is the number of times the winning PAN presidential candidate visited a district in 2006 leading up to the 2009 federal elections. A count here means that is a number of times, fixed and a discrete variable. \newline

\noindent Our main predictor of interest is whether the district was highly contested or  whether it was not (the PAN or their opponents have electoral security) in the  previous federal elections during 2000 (\texttt{competitive.district}), which is binary (1 = close/swing district, 0 = ``safe seat''). We also include  \texttt{marginality.06} (a measure of poverty) and \texttt{PAN.governor.06}  (a dummy for whether the state has a PAN-affiliated governor) as additional  control variables.\newline

\noindent When I check the structure of the data using \texttt{str()}, I confirm that 
\texttt{PAN.visits.06} is an integer. The median is 0 and the mean is 0.09, which 
implies that we have many zeros in our dataset. This makes sense, because a 
candidate might have never visited a district, which contains equal relevance in our analysis as when the candidate has visited one time. In other words, an outcome of 0 has equal meaning.\newline

\noindent Next, I check whether the mean and variance are approximately equal, as this is the first red flag for Poisson. The mean is 0.09 and the variance is 0.64, meaning they are not equal. This means there is overdispersion in our data set, impying  the presence of greater variability in our data than expected for this model (Poisson).\newline

\noindent Because the outcome is a count variable, OLS is not appropriate here.  OLS would predict negative values and we want zeros or positives. Furthermore, as the number of visits increases, the dispersion also increases. Therefore I use a Poisson regression with a log link.\newline

\tiny \begin{BVerbatim}
	mexico_q2 <- na.omit(mexico_elections[, c(
	"PAN.visits.06", "competitive.district",
	"marginality.06", "PAN.governor.06"
	)])
	
	mean(mexico_elections$PAN.visits.06)
	var(mexico_elections$PAN.visits.06)
	
	pois_mod <- glm(
	PAN.visits.06 ~ competitive.district + marginality.06 + PAN.governor.06,
	data   = mexico_q2,
	family = poisson(link = "log")
	)
	
	summary(pois_mod)
	
	comp_z <- coef(summary(pois_mod))["competitive.district", "z value"]
	comp_p <- coef(summary(pois_mod))["competitive.district", "Pr(>|z|)"]
	comp_z
	comp_p
	
	exp(coef(pois_mod))
	
	new_case <- data.frame(
	competitive.district = 1,
	marginality.06       = 0,
	PAN.governor.06      = 1
	)
	predict(pois_mod, newdata = new_case, type = "response")
\end{BVerbatim}
\normalsize \newline

\begin{table}[H]
	\centering
	\caption{Poisson Regression Model: PAN Candidate Visits in 2006}
	\label{tab:poisson}
	\begin{tabular}{lcccc}
		\hline\hline
		& Estimate        & Std. Error & z value & IRR    \\
		\hline
		(Intercept)          & -3.8102$^{***}$ & 0.2221     & -17.156 & 0.0221 \\
		competitive.district & -0.0814         & 0.1707     & -0.477  & 0.9219 \\
		marginality.06       & -2.0801$^{***}$ & 0.1173     & -17.728 & 0.1249 \\
		PAN.governor.06      & -0.3116$^{*}$   & 0.1667     & -1.869  & 0.7323 \\
		\hline
		\multicolumn{5}{l}{\textit{Model fit}} \\
		Null Deviance        & \multicolumn{4}{c}{1473.87} \\
		Residual Deviance    & \multicolumn{4}{c}{991.25}  \\
		AIC                  & \multicolumn{4}{c}{1299.2}  \\
		\hline\hline
		\multicolumn{5}{l}{\footnotesize $^{*}$p$<$0.1; $^{**}$p$<$0.05; $^{***}$p$<$0.01} \\
		\multicolumn{5}{l}{\footnotesize IRR = Incidence Rate Ratio = 
			exp(coefficient).} \\
	\end{tabular}
\end{table}

\subsubsection{A: Do PAN candidates visit swing districts more?}

\noindent To answer this question, I look at the coefficient for \texttt{competitive.district}. The z-statistic is -0.477 and the p-value is 0.6336394. This means we fail to reject the null hypothesis that PAN presidential candidates visit swing districts more often than safe seats. There is no siginificant evidence for this effect.\newline

\subsubsection{B: Interpret marginality.06 and PAN.governor.06}

 \noindent The coefficients are difficult to directly interpret, therefore I take 
\texttt{exp()} and transform them into Incidence Rate Ratios (IRRs). An IRR tells 
us by what multiplicative factor the expected number of visits changes for a 
one-unit increase in the predictor, holding everything else constant. An IRR $>$ 1  means more visits, an IRR $<$ 1 means fewer visits.\newline

\noindent For \texttt{marginality.06}, the coefficient is -2.08 and highly significant (p $<$ 0.001), with an IRR of 0.12. The expected number of visits 
decreases by a multiplicative factor of 0.12 for each one-unit increase in the marginality index. Put simply, districts with higher poverty levels receive about 88\% fewer visits compared to less poor districts, holding everything else 
constant.\newline

\noindent For \texttt{PAN.governor.06}, the coefficient is -0.31 and marginally 
significant (p $<$ 0.1), with an IRR of 0.73. This means that districts in states 
with a PAN-affiliated governor receive about 27\% fewer visits compared to 
districts without a PAN governor, holding everything else constant. However, given the marginal significance, this should be interpreted with caution.\newline

\subsubsection{C: Predicted number of visits}

\noindent Using \texttt{predict()} with \texttt{type = "response"}, which gives 
the expected count, I predict the mean number of visits for a hypothetical district  that was competitive (\texttt{competitive.district} = 1), had an average poverty level (\texttt{marginality.06} = 0), and had a PAN governor 
(\texttt{PAN.governor.06} = 1). The model predicts \textbf{0.015} visits. This means that for a competitive district with average poverty and a PAN-affiliated 
governor, the model predicts 0.015 visits (on average, because it is calculating mean number of visits). This is an extremely low number, which makes sense given that we have a lot of zeros in our dataset. In fact most districts received zero visits, and the mean across all  districts was only 0.09.\newline

\end{document}
